\chapter{Обзор научно-технических источников информации}\label{ch:1}

Задача поиска повреждений автомобилей совсем не нова - количество научных источников по теме автоматизации этого процесса с помощью искусственного интеллекта достаточно велико. Большинство из них описывают высокую коммерческую значимость проблемы - одни упоминают страховые компании, другие - продажи подержанных автомобилей, третьи - ремонт в автосервисах. Так, в одном из источников указано, что страховые компании США ежегодно теряют порядка 18 миллиардов долларов из-за ручной обработки обращений \cite{src3}. Эти цифры и развитие методов машинного обучения требуют внедрять решения, направленные на ускорение и упрощение оценки с помощью компьютерного зрения. В рамках рассмотренных статей исследуются методы автоматизации процесса визуальной оценки повреждений автомобилей. 

В данном разделе отражены имеющиеся подходы к реализации задачи поиска и классификации повреждений по фотографии, содержащей автомобиль. Наиболее часто применяемый подход - использование разного рода архитектур свёрточных нейронных сетей (CNN).

В масштабном анализе и сводке по подходам машинного обучения к решению задачи обнаружения повреждений автомобилей, проведённом в 2025 году, подробно рассказано о применяемых методах. Их можно разделить на 2 группы - традиционные методы машинного обучения и методы, основанные на глубоком обучении \cite{src1}. Для каждой категории подходы дополнительно группируются по трем ключевым применениям: классификация, обнаружение и сегментация.

\begin{figure}[h]
\centering
\includegraphics[scale=2]{Общая схема подходов CNN}
\caption{Обзор системы обнаружения повреждений транспортных средств \cite{src1}}
\end{figure}


Классификация - для определения типа повреждения/детали, обнаружение (детекция) - для определения местоположения, а сегментация - для выделения точных областей, обеспечения точности на уровне пикселей.


Проведённый глубокий поиск статей и источников показал, что возможности традиционных методов ограничены только простыми задачами классификации, такими как различения повреждённых и неповреждённых транспортных средств или определения базовых типов повреждений. По более сложным задачам научных статей обнаружено не было. Это объясняется тем, что определение границ поврежденных областей - трудная задача для традиционных методов по причине их зависимости от низкоуровневых объектов и ограниченного пространственного контекста. \cite{src1}. В одной из статей также ссылаются на источник, в котором проводилось исследование, показавшее, что в задаче обнаружения мелких повреждений в стальных канатах свёрточные нейронные сети значительно превосходят другие методы машинного обучения, такие как kNN, SVM, Random Forest и Logistic Regression. \cite{src6}


В подавляющем большинстве рассмотренных современных работ по распознаванию и классификации повреждений транспортных средств используются CNN-архитектуры для задач классификации, детекции и сегментации. Рассматриваемые методы нацелены на быстрый и информативный анализ изображений - именно экономия времени и ресурсов при сохранении требуемого уровня точности является основной целью исследований.

\section{Свёрточные нейронные сети в анализе изображений}\label{ch:1.1}

Глубокое обучение, в частности свёрточные нейронные сети (CNN), продемонстрировало впечатляющие результаты в широком спектре задач визуального распознавания, от обнаружения объектов до семантической сегментации \cite{src7} . 

Основная идея CNN заключается в последовательном выделении признаков с изображения, начиная с простых граней, текстур на первых уровнях и заканчивая сложными формами в более глубоких слоях. 

Типовая CNN состоит из последовательности сверточных слоёв, выделяющих локальные признаки изображения, и слоёв объединения (пуллинговых слоёв) для уменьшения размерности и устойчивости к смещению. После свёрток обычно идут полносвязные слои (или сверточные «головы») для предсказания классов или координат объектов. \cite{src14}

\begin{figure}[h]
\centering
\includegraphics[scale=0.65]{Typical_CNN}
\caption{Концепция архитектуры CNN \cite{src14}}
\end{figure}


Такая архитектура позволяет выявлять сложные визуальные закономерности на изображении. Как отмечено в одном из источников, CNN особенно хороши в распознавании краёв, форм, текстур на изображениях, что делает их отличными помощниками в классификации изображений. \cite{src3}

Современные модели с CNN архитектурой можно разбить на две большие группы - одноэтапные и двухэтапные. 

Представителями двухэтапных методов является серия R-CNN. (Mask R-CNN, Faster R-CNN) Этот метод генерирует множество областей-кандидатов, которые могут содержать объекты, при первой обработке изображения. Затем для каждой области-кандидата производится повторная обработка через CNN. Это обеспечивает высокую точность предсказания, однако есть и главный недостаток: подход достаточно медленный и требует больших вычислительных ресурсов.
Поэтому был разработан алгоритм YOLO (You Only Look Once), впервые предложенный в 2016 году - он обрабатывает всё изображение "за один проход", что ускоряет обнаружение объектов. \cite{src2}

Серия YOLO получила широкую популярность благодаря балансу скорости и точности. Начиная с YOLOv5 и заканчивая YOLOv8, а в последнее время и YOLOv11, модели претерпели значительные архитектурные переработки, такие как улучшение базовых компонентов, подходов к агрегированию признаков и методов обучения, что в совокупности повысило эффективность обнаружения. Несмотря на улучшения и прочие достижения, в сфере обнаружения повреждений транспортных средств по-прежнему существуют проблемы. \cite{src7}

Рассмотрим более подробно существующие решения задачи оценки повреждений автомобилей.

\section{Разделение задачи на подзадачи: детали, повреждения и оценка стоимости}\label{ch:1.2}


Задача оценки повреждений достаточно масштабная. И нередко её разбивают на несколько взаимосвязанных этапов, таких как определение повреждений автомобиля и оценка стоимости. Так, в статье "Vehicle damage severity estimation for insurance operations using inthe-wild mobile images" описывается подход пайплайна подзадач: они обнаруживают автомобиль, затем проводят семантическую сегментацию всех его деталей и всех повреждённых участков, рассматривая и тот, и другой этап как задачу инстанс-сегментации с применением Mask R-CNN. После сегментации всех кузовных частей и всех повреждённых частей вычисляется относительная площадь повреждений по каждой детали, и затем передаётся в модель окончательной оценки стоимости ремонта. В работе шаг окончательной оценки стоимости и прочие шаги также выделяются как отдельные части пайплайна.  \cite{src9}

\begin{figure}[h]
\centering
\includegraphics[scale=0.5]{Пайплайн}
\caption{Обзор пайплайна процесса оценки ущерба \cite{src9}}
\end{figure}

Такой подход имеет несколько преимуществ:

\begin{itemize}
   \item Это позволяет последовательно решать проблему, поэтапно выполняя действия по получению нужного результата от анализа изображения.
   \item Объекты разных типов (детали и повреждения) могут требовать разных подходов для необходимой точности обнаружения, а предварительное разбитие на подзадачи сегментации позволяет снять часть фона и точнее анализировать повреждения.
   \item Разделение на подзадачи позволяет выполнить его с помощью нескольких моделей (например, одна модель для сегментации деталей, а вторая - для повреждений). Это позволяет сначала обучить модель находить и классифицировать детали, а затем распознавать повреждения на уже вырезанных деталях
   \item Пайплайн из подзадач позволяет видеть промежуточные результаты и анализировать их. Ошибки и прочие проблемы на каждой из фаз общего процесса можно корректировать отдельно
\end{itemize}

Продолжая тему разделения ответственности решения на несколько этапов, рассмотрим подходы по детекции и классификации как деталей, так и повреждений.


\section{Подходы к распознаванию деталей автомобиля}\label{ch:1.3}


При анализе изображения на предмет обнаружения деталей автомобиля применяются как сегментационные, так и детекторные методы.
В ряде работ реализована инстанс-сегментация деталей, позволяющая выделять деталь отдельной маской. Существует множество алгоритмов сегментации, каждый из которых имеет свои преимущества и недостатки. 
В статье "Evaluation of deep learning algorithms for semantic segmentation of car parts" проведено сравнение пяти алгоритмов сегментации - Mask R-CNN, Hybrid Task Cascade (HTC), Composite Backbone Network (CBNet), Path Aggregation Network (PANet) и Global Context Network (GCNet)), в котором best-performing оказался HTC с ResNet-50 по метрике mAP. Авторы работы также оценивали устойчивость модели к зашумлённым данным (в частности, при плохих погодных условиях - туман, снегопад) и GCNet показала лучший результат по метрике mPC. \cite{src10}

Параллельно с сегментацией, детали часто определяют стандартной детекцией. Так, в работе "Optimized car parts detection with advanced feature fusion and attention modules" предложена улучшенная YOLO-архитектура для детекции мелких частей автомобиля. \cite{src11}

\subsection{Mask R-CNN}\label{ch:1.3.1}


Алгоритм сегментации объектов Mask R-CNN был разработан на основе Faster R-CNN. Его новизна заключается в том, что Faster R-CNN мог только определять местоположение объекта на изображении и распознавать его, а Mask R-CNN также может производить и сегментацию объектов.

Mask R-CNN состоит из двух основных частей:
\begin{itemize}
   \item Базовой модели (Backbone), которая извлекает признаки из изображения с помощью остаточной нейронной сети (ResNet), состоящей из 50–101 слоя свёрточной нейронной сети, в сочетании с сетью Feature Pyramid Network (FPN)
   \item "Головы" (Head), которая создаёт ограничивающую рамку (Bounding Box) вокруг интересующей области (ROI) и определяет тип объекта в этой рамке
\end{itemize}

Ещё одна особенность, отличающая Mask R-CNN от Faster R-CNN - создание маски для каждой области интереса (ROI).  \cite{src10}

\begin{figure}[h]
\centering
\includegraphics[scale=0.5]{Mask R-CNN}
\caption{Архитектура Mask R-CNN \cite{src10}}
\end{figure}


\subsection{Global Context Network (GCNet)}\label{ch:1.3.2}

Архитектура Global Context Network (GCNet) имеет структуру, схожую с Mask R-CNN, которая показана на рисунке 4. Однако особенностью GCNet является расширение Backbone блоком глобального контекста, схема которого отражена на рисунке 5.

\begin{figure}[h]
\centering
\includegraphics[scale=0.5]{GCNet}
\caption{Блок глобального контекста (Global Context). Карта признаков имеет размер C × H × W — количество каналов C, высота H и ширина W. × обозначает матричное умножение, а + — поэлементное сложение с трансляцией. r — коэффициент узкого места, а C/r — размерность скрытого представления узкого места \cite{src10}}
\end{figure}

Преимущество такого блока в том, что в его состав входит Non-Local Network (NLNet), помогающая решить проблему дальнодействующих зависимостей. \cite{src10}

\subsection{Hybrid Task Cascade (HTC)}\label{ch:1.3.3}

Алгоритм Hybrid Task Cascade (HTC) был создан с целью повышения эффективности инстанс-сегментации. 

Особенности:
\begin{itemize}
    \item В архитектуре блоки предсказания ограничивающих рамок (bounding box'ов), масок и извлечения ROI организованы в каскадную структуру, как показано на рисунке 6
    \item Процессы предсказания ограничивающих рамок и масок выполняются параллельно и последовательно на нескольких стадиях
    \item Многостадийная ветвь предсказания масок - маска, полученная на одном этапе, используется при генерации маски на следующем.
    \item Модель учитывает контекст при формировании каждой маски
\end{itemize}

\begin{figure}[h]
\centering
\includegraphics[scale=0.5]{HTC}
\caption{Архитектура HTC \cite{src10}}
\end{figure}

Совокупность этих улучшений существенно улучшает передачу информации между задачами и повышает итоговую точность. \cite{src10}

\subsection{Усовершенствованная архитектура на базе YOLO}\label{ch:1.3.4}

Важные особенности архитектуры:
\begin{itemize}
    \item Внедрён C2f (Cross Stage Partial style block) блок - помогает сохранять низкоуровневые признаки
    \item Внедрён модифицированный C2fCIB (C2f + Channel Interaction Block) блок - усиливает межканальные зависимости
    \item PSA (Pyramid Split Attention Module) - повышает выделение мелких/схожих деталей
\end{itemize}

\begin{figure}[h]
\centering
\includegraphics[scale=0.5]{improved_yolo}
\caption{Архитектура улучшенной YOLO \cite{src11}}
\end{figure}

Однако подход имеет и недостатки:
\begin{itemize}
    \item Очень мелкие части (<1\% изображения) всё так же остаются проблемой
    \item Модель стала заметно тяжелее по параметрам (3.2M параметров у YOLOv8 против 74.6M у предложенной модели) и требует куда больше вычислений. \cite{src11}
\end{itemize}

Результаты сравнения с другими подходами показали, что предложенные архитектурные улучшения дают существенный рост по метрикам Recall и mAP в задаче детекции деталей автомобиля. Авторский вариант YOLO превосходит некоторые классические двухэтапные R-CNN подходы, такие как Faster R-CNN и SSD, по метрикам Precision, Recall, mAP на рассматриваемом датасете car-parts.

\section{Подходы к распознаванию повреждений автомобиля}\label{ch:1.4}

Методы распознавания повреждений автомобиля, в отличие от задачи детекции автомобильных деталей, требуют более тщательного и точного анализа изображений. По мнению авторов статьи "Vehicle Damage Severity Estimation for Insurance Operations Using In-The-Wild Mobile Images", задача обнаружения повреждений скорее является задачей сегментации, а не задачей детекции, и многие современные методы отдают предпочтение именно сегментации. Сегментационные маски позволяют точнее описывать размытые очертания и мелкие дефекты, чем ограничивающие рамки (bounding box'ы). \cite{src9}

В статье "Automatic damaged vehicle estimator using enhanced deep learning algorithm" используют улучшенный Mask R-CNN с предварительно обученными моделями Inception-ResNetV2, VGG16, VGG19 для локализации и классификации повреждений. Комбинация Mask R-CNN с Inception-ResNetV2 показала лучшие результаты \cite{src4}

Популярны так же и решения с улучшенными YOLO-архитектурами специально для повреждений. 

Один из примером улучшенного YOLO - HL-YOLO. Это модифицированный YOLOv11, использующий гетерогенные свёртки и специальные механизмы внимания.

Особенности архитектуры:
\begin{itemize}
    \item Используются механизмы внимания для более точного выявления мелких и нерегулярных повреждений автомобиля, таких как царапины и вмятины.
    \item Применяется Large Separable Kernel Attention (LSKA), позволяющий учитывать пространственный контекст на больших расстояниях и повышать чувствительность к слабовыраженным, мелким дефектам.
    \item Дополнительным преимуществом HL-YOLO является использование модуля C3k2-HetConv, который включает гетерогенные свёртки и частичное разделение признаков. Это помогает одновременно и снизить вычислительные затраты, и повысить чувствительность к мелким повреждениям \cite{src7}
\end{itemize}

\begin{figure}[h]
\centering
\includegraphics[scale=1]{hl_yolo}
\caption{Архитектура HL-YOLO \cite{src7}}
\end{figure}

Также в источнике "Small Dents, Big Impact: A Dataset and Deep Learning Approach for Vehicle Dent Detection" был проведён следующий эксперимент: собран датасет мелких вмятин и обучены YOLOv8-модели. Модель YOLOv8m-t42 достигла точности ~0.86 и полноты ~0.84 при обнаружении даже микроскопических вмятин. Полученные результаты указывают на то, что одностадийные детекторы тоже могут быть эффективны для этой узкой задачи нахождения повреждений. \cite{src12}

Проведённые сравнения моделей говорят о том, что выбор зависит от требований к точности и скорости. В одном из сравнений YOLO и Faster R-CNN  показано, что YOLO показывают результат по метрикам Recall и mAP немного хуже, чем Faster R-CNN, однако требуют меньше ресурсов (GFLOPs) и намного быстрее (FPS). \cite{src7}

\begin{figure}[h]
\centering
\includegraphics[scale=0.5]{yolo_vs_r_cnn}
\caption{Точность и эффективность различных моделей обнаружения \cite{src7}}
\end{figure}

Исследование показывает, что детекторы YOLO дают баланс скорости и точности, в то время как инстанс-сегментаторы типа Mask R-CNN могут обеспечивать более высокую точность (маски повреждений), но с использованием больших ресурсов.


\begin{table}[ht]
\caption{Архитектурные подходы к решению задачи оценки повреждений}
\label{table:simple_example}
\begin{tabularx}{\textwidth}{|>{\raggedright\arraybackslash}p{4.5cm}
                          |>{\raggedright\arraybackslash}X
                          |>{\raggedright\arraybackslash}X|}
\hline
\textbf{Описание архитектурного подхода} & \textbf{Минусы} & \textbf{Плюсы} \\
\hline
Общая модель и для деталей, и для повреждений
&
Сложнее обучать и интерпретировать результаты
&
Минимальное кол-во этапов обработки
\\
\hline
Пайплайн подзадач (детали и повреждения как отдельные шаги)
&
Увеличение числа этапов
&
Улучшенная управляемость, точность, возможность оптимизации каждого этапа
\\
\hline
\end{tabularx}
\end{table}



\begin{table}[ht]
\caption{Подходы(модели) для сегментации деталей автомобиля}
\label{table:simple_example}
\begin{tabularx}{\textwidth}{|>{\raggedright\arraybackslash}p{4.5cm}
                          |>{\raggedright\arraybackslash}X
                          |>{\raggedright\arraybackslash}X|}
\hline
\textbf{Описание подхода / модели} & \textbf{Минусы} & \textbf{Плюсы} \\
\hline
Mask R-CNN
&
Высокая сложность, скорость инференса ниже
&
Очень высокая точность масок, чёткое разделение деталей
\\
\hline
Hybrid Task Cascade (HTC)
&
Более сложная архитектура и обучение
&
Достаточно высокая точность - за счёт каскадного уточнения масок
\\
\hline
Global Context Network (GCNet)
&
Увеличение сложности модели
&
Устойчивость к нестандартным условиям съёмки, шумам
\\
\hline
YOLO-архитектуры
&
Могут выдавать недостаточно точные маски деталей
&
Высокая скорость обучения и инференса
\\
\hline
\end{tabularx}
\end{table}

\begin{table}[ht]
\caption{Подходы(модели) для сегментации повреждений автомобиля}
\label{table:simple_example}
\begin{tabularx}{\textwidth}{|>{\raggedright\arraybackslash}p{4.5cm}
                          |>{\raggedright\arraybackslash}X
                          |>{\raggedright\arraybackslash}X|}
\hline
\textbf{Описание модели} & \textbf{Минусы} & \textbf{Плюсы} \\
\hline
Mask R-CNN 
&
Требует точной разметки, много ресурсов, дольше инференс
&
Точно выделяет повреждения, лучше локализация
\\
\hline
YOLOv8
&
Может выдавать плохие результаты на сложных повреждениях
&
Баланс скорости и точности результатов
\\
\hline
Улучшенные YOLO-архитектуры
&
Недостаточно точные границы повреждений
&
Повышенная чувствительность к мелким дефектам
\\
\hline
\end{tabularx}
\end{table}

\begin{table}[ht]
\caption{Сравнение двухэтапных и одноэтапных подходов}
\label{table:simple_example}
\begin{tabularx}{\textwidth}{|>{\raggedright\arraybackslash}p{4.5cm}
                          |>{\raggedright\arraybackslash}X
                          |>{\raggedright\arraybackslash}X|}
\hline
\textbf{Подход} & \textbf{Минусы} & \textbf{Плюсы} \\
\hline
Двухэтапные (Mask R-CNN, HTC и др.)
&
Медленные, требовательны к ресурсам и точной разметке
&
Максимальная точность сегментации
\\
\hline
Одноэтапные (YOLO и модификации)
&
Менее точные маски
&
Высокая скорость и более простое обучение
\\
\hline
\end{tabularx}
\end{table}
